\section{The primary structure of quartic hyperplane multiplication}
\label{sec:quartic-primary}
Put $T_m=m(m+1)/2$. The order of the maximal-minor ideal at an
evaluation is governed by $T_m$, not merely by the corank $m$.
For quartic binary forms this order determines the entire ideal.
Throughout this section the ideals of the evaluation curve and the
secant variety are reduced ideal sheaves.

\begin{theorem}[Quartic primary structure in every degree]
\label{thm:quartic-primary}
Let $n=4$ and $m\ge2$. On $\Pj(V_4^*)$ one has
\begin{equation}\label{eq:quartic-primary-ideal}
 \mathcal J_m
 =\mathcal I_{S_4}\mathcal I_{C_4}^{T_m-2}
 =\mathcal I_{S_4}\cap\mathcal I_{C_4}^{T_m}.
\end{equation}
The second equality is an irredundant primary decomposition of ideal
sheaves: its associated subvarieties are precisely $S_4$ and $C_4$.
At every $x\in C_4$ there are formal coordinates $s,A,B,C$ such that
\begin{equation}\label{eq:quartic-complete-germ}
 \widehat{\mathcal O}_{D_m,x}
 \simeq
 \C[[s,A,B,C]]\big/
       \big((AC-B^2)(A,B,C)^{T_m-2}\big),
 \qquad \mathcal I_{C_4}=(A,B,C).
\end{equation}
Consequently the nilpotency index of the nilradical at $x$ is exactly
$\lceil T_m/2\rceil$. The transverse nilradical has length
$\binom{T_m}{3}$ at the generic point of $C_4$.
\end{theorem}

The assertion is about the actual maximal-minor ideal, not only its
support or its associated graded ideal. The passage from the lowest
terms to the complete ideal will use a uniform spanning statement
for minors. We first establish the required filtration in all binary
degrees.

\subsection{A filtered presentation along the evaluation curve}
Work on $h_0=1$, with $h_j$ the value of the functional on $t^j$.
The translation $t\mapsto t-h_1$ gives the polynomial coordinates
\[
 s=h_1,\qquad
 a_j=\sum_{i=0}^j\binom ji(-s)^{j-i}h_i\quad(2\le j\le n).
\]
Thus $a_0=1$, $a_1=0$, and the evaluation curve is defined by
$\mathfrak p=(a_2,\ldots,a_n)$. The inverse change is the opposite
translation. Since multiplication commutes with translation, its
matrix is equivalent, by invertible changes of source and target
over $\C[s,a_2,\ldots,a_n]$, to one independent of $s$. We may therefore
use the basis
\begin{equation}\label{eq:centered-hyperplane-basis}
 f_1=t,\qquad f_j=t^j-a_j\quad(2\le j\le n)
\end{equation}
for $U$, with homogeneous powers of the second binary variable
understood. No affine-degree truncation is imposed on a product:
its target is the full space of polynomials of degree at most $mn$.

\begin{lemma}[Row filtration]
\label{lem:triangular-row-filtration}
For every $n\ge4$ and $m\ge2$, elementary operations present
$\operatorname{coker}\mu_m$ by a matrix with $m$ rows, indexed by
$k=0,\ldots,m-1$, whose row $k$ has entries in
$\mathfrak p^{m-k}$. Consequently
\begin{equation}\label{eq:triangular-vanishing-order}
 \mathcal J_m\subset\mathcal I_{C_n}^{T_m}.
\end{equation}
At each point of $C_n$ the order is exactly $T_m$.
\end{lemma}
\begin{proof}
Every integer from $m$ to $mn$ is the sum of $m$ integers in
$[1,n]$. Choose one product of the $f_j$ for each such sum. Their
coefficient matrix in rows $t^m,\ldots,t^{mn}$ is triangular with
diagonal entries one. It can be eliminated over the polynomial
coefficient ring, without localization. The rows left over are
$t^0,\ldots,t^{m-1}$.

In a product of $m$ factors, a term contributing to $t^k$ uses at
most $k$ nonconstant factors. Its coefficient therefore contains at
least $m-k$ of the parameters $a_j$. Every pivot column satisfies
the same inequality. Subtracting polynomial multiples of pivot
columns preserves it, proving the row assertion and the inclusion.

For completeness, an explicit nonzero initial minor proves equality
of orders. Write $A=a_2$, $B=a_3$, $C=a_4$, give these variables
weights $1,2,3$, and index the residual rows by $q=m-k$.
If $P$ is a monomial of ordinary degree $q$ and weight $w$, it
corresponds to a source product with $m-q$ factors $f_1$ and with
leading $t$-degree $m+w$. After absorbing the row sign $(-1)^q$,
its row-$q$ initial term is $P$. Two such products of the same
weight have cancelling top terms; all subsequent pivot coefficients
lie in $\mathfrak p$. Their residual difference therefore has initial
vector $P e_q-Q e_r$, where $Q$ has degree $r$ and the same weight.
Here and below ``initial'' in row $q$ means its part of ordinary
degree $q$.

Take the $m-1$ differences with initial vectors
\[
 B A^{q-1}e_q-A^{q+1}e_{q+1},\qquad 1\le q<m,
\]
and the final difference with initial vector
$A^{m-2}(AC-B^2)e_m$. Their determinant has initial term, up to sign,
\[
 A^{(m-2)(m+1)/2}B^{m-1}(AC-B^2),
\]
which is nonzero and has degree $T_m$. The determinant belongs to
$\mathcal J_m$ by multilinearity, since the columns used are residual
images of source-column differences. This construction uses only
$f_1,\ldots,f_4$ and is valid for every $n\ge4$. Translation and
homogeneous changes of coordinates cover $C_n$.
\end{proof}

\begin{corollary}[All-degree embedded support and quadratic growth]
\label{cor:all-degree-embedded-support}
For every $n\ge4$ and $m\ge2$,
\[
 \operatorname{Ass}(\mathcal O_{D_m})
       =\{\eta_{S_n},\eta_{C_n}\}.
\]
At every evaluation point,
$\mathcal N_m^j\ne0$ for $2j<T_m$, and hence the nilpotency index
is at least $\lceil m(m+1)/4\rceil$.
\end{corollary}
\begin{proof}
In centered coordinates the secant equation obtained from the first
$3$ by $3$ Hankel minor is
$a_2a_4-a_3^2-a_2^3$. Its $\mathfrak p$-order is two. Its $j$-th
power cannot lie in $\mathcal J_m$ when $2j<T_m$, by
Lemma~\ref{lem:triangular-row-filtration}. Since $T_m\ge3$, the
nilradical is nonzero already for $m=2$. Smoothness off $C_n$ and
Lemma~\ref{lem:equivariant-associated-curve} give the associated-point
assertion exactly as in Theorem~\ref{thm:higher-hyperplane-global}.
\end{proof}

\subsection{A uniform spanning lemma for the lowest minors}
Let $R=\C[A,B,C]$, $\mathfrak p=(A,B,C)$, and $F=AC-B^2$.
In addition to ordinary degree, use the weights $1,2,3$ above.
Let $E_m$ be the ideal of $m$-minors of the finite matrix whose
columns are all vectors
\begin{equation}\label{eq:weighted-relation-columns}
 P e_q-Q e_r,\qquad
 1\le q,r\le m,\quad
 \deg P=q,\quad \deg Q=r,\quad
 \operatorname{wt}P=\operatorname{wt}Q.
\end{equation}
Repeated or zero columns may be omitted. Terms with $q=r$ are
allowed; they give one-row columns. Every entry in row $q$ has
ordinary degree $q$.

\begin{lemma}[Weighted-relation minors]
\label{lem:weighted-relation-minors}
For every $m\ge2$,
\begin{equation}\label{eq:weighted-minor-identity}
 E_m=F\mathfrak p^{T_m-2}.
\end{equation}
\end{lemma}
\begin{proof}
Each maximal minor is homogeneous of degree $T_m$. Substituting
$(A,B,C)=(u,uv,uv^2)$ makes the matrix have a nonzero left kernel:
row $q$ is multiplied by $(v/u)^q$, and the two monomials in each
column then both have value $v^{\operatorname{wt}P}$. This computation
is over $\C(u,v)$. Since the kernel of
$R\to\C[u,v]$ is $(F)$, every minor is divisible by $F$. Thus
$E_m\subset F\mathfrak p^{T_m-2}$.

For $m=2$ the three columns
$Be_1-A^2e_2$, $Ce_1-ABe_2$, and $Fe_2$ have minors
$AF,BF,CF$, up to signs. This proves the base case.
Assume the assertion for $m$, and set $K=T_{m+1}-2$.
The matrix for $m$ embeds in the first $m$ rows of the matrix for
$m+1$, with zero last row. If $P$ has degree $m+1$ and weight at
most $3m$, there is a monomial $Q$ of some degree $r\le m$ with
the same weight: the intervals $[r,3r]$, $1\le r\le m$, cover
$[1,3m]$. Appending $Pe_{m+1}-Qe_r$ to any old maximal minor shows
that
\begin{equation}\label{eq:weighted-induction-extension}
 P E_m\subset E_{m+1}.
\end{equation}

For a degree-$K$ monomial $M$, put
$\delta(M)=3K-\operatorname{wt}M$; thus an $A$ contributes two to
$\delta$ and a $B$ contributes one. If $\delta(M)\ge3$, one can
select at most three factors of $M$ having total deficit at least
three, and complete them to a degree-$m+1$ factor $P$ of $M$.
The remaining degree is $T_m-2$. The weight of $P$ is at most
$3m$, so \eqref{eq:weighted-induction-extension} and the induction
hypothesis give $FM\in E_{m+1}$.

Only $C^K$, $BC^{K-1}$, $AC^{K-1}$ and $B^2C^{K-2}$ remain.
Use the $m$ columns
\[
 C^qe_q-ABC^{q-1}e_{q+1}\quad(1\le q\le m)
\]
and a final one-row column $FH e_{m+1}$, where $H$ is any monomial
of degree $m-1$. Their determinant is $FC^{T_m}H$, up to sign.
For $m\ge3$, choosing
\[
 H=C^{m-1},\quad BC^{m-2},\quad AC^{m-2},\quad B^2C^{m-3}
\]
gives the four missing generators. When $m=2$, the first three
choices remain available. The columns
$Ce_1-ABe_2$, $Fe_2$, and $CFe_3$ additionally give $C^2F^2$.
The identity
$AC^3F-C^2F^2=B^2C^2F$ supplies the last generator.
This proves the induction, and hence the equality of ideals.
\end{proof}

\subsection{From lowest terms to the complete primary ideal}
\begin{proof}[Proof of Theorem~\ref{thm:quartic-primary}]
Use the centered coordinates above with $n=4$, and make the
triangular change
\[
 A=a_2,\qquad B=a_3,\qquad C=a_4-a_2^2.
\]
The secant equation becomes $F=AC-B^2$. Its irreducibility and the
three-dimensionality of the secant variety show that
$\mathcal I_{S_4}=(F)$ on this chart. The ideal of $C_4$ is
$\mathfrak p=(A,B,C)$.

Write $J$ for the local maximal-minor ideal. The support theorem
and Lemma~\ref{lem:triangular-row-filtration} imply
\[
 J\subset (F)\cap\mathfrak p^{T_m}
       =F\mathfrak p^{T_m-2}=:I.
\]
The last equality follows from the homogeneous quadratic equation
$F$: multiplication by $F$ adds exactly two to $\mathfrak p$-order,
also over the coefficient ring $\C[[s]]$.

For every column \eqref{eq:weighted-relation-columns}, replace its
monomial $P$ by the product of $f_2,f_3,f_4$ specified by its
exponents, and insert $m-q$ factors $f_1$. Do the same for $Q$.
Both source products have degree $m+\operatorname{wt}P$ in $t$;
their difference has cancelling top coefficient. As proved in the
row-filtration lemma, its residual initial column is exactly
\eqref{eq:weighted-relation-columns}, after a fixed change of row
signs. The change $a_4=C+A^2$ does not change these lowest terms.
Thus the degree-$T_m$ initial forms of elements of $J$ include all
maximal minors defining $E_m$. Lemma~\ref{lem:weighted-relation-minors}
says that they span the whole degree-$T_m$ part of $I$.

It follows that $J+\mathfrak p I=I$. Indeed the difference between
a chosen lift in $J$ and its homogeneous generator lies in
$I\cap\mathfrak p^{T_m+1}=\mathfrak p I$. Nakayama's lemma,
applied to the finite module $I/J$ in the completed local ring,
gives $J=I$. Completion is faithfully flat, so this is an equality
of local ideals as well. Away from $C_4$, the equality follows from
Theorem~\ref{thm:higher-hyperplane-global}. The ideal-sheaf identities
\eqref{eq:quartic-primary-ideal} and the completed germ follow.

The ideal $(F)$ is prime, and $\mathfrak p^{T_m}$ is
$\mathfrak p$-primary in these coordinates. Neither component can
be omitted. This proves the asserted primary decomposition. A power
$F^j$ belongs to $F\mathfrak p^{T_m-2}$ exactly when
$2j\ge T_m$, proving the exact nilpotency index. Finally the
nilradical, as a module, is
$F R/F\mathfrak p^{T_m-2}\simeq R/\mathfrak p^{T_m-2}$.
Over $\C(C_4)$ its length is the number of monomials of degree
less than $T_m-2$ in three variables, namely $\binom{T_m}{3}$.
\end{proof}

\begin{corollary}[Infinitesimal layers and normal cone]
\label{cor:quartic-infinitesimal-layers}
Let $\mathcal N=\mathcal I_{S_4}/\mathcal J_m$ and let
$\mathcal E=\mathcal I_{C_4}/\mathcal I_{C_4}^2$ be the conormal
bundle. For $0\le j<T_m-2$,
\begin{equation}\label{eq:quartic-conormal-layers}
 \frac{\mathcal I_{C_4}^j\mathcal N}
      {\mathcal I_{C_4}^{j+1}\mathcal N}
 \simeq
 \mathcal O_{C_4}(-3)\otimes\Sym^j\mathcal E.
\end{equation}
Here $\mathcal O_{C_4}(-3)$ is the restriction of
$\mathcal O_{\Pj^4}(-3)$, hence has degree $-12$ on $\Pj^1$.
For $j\ge1$ and $2j<T_m$, the nilradical powers and their quotients
are
\begin{align}
 \mathcal N^j
 &\simeq\mathcal I_{S_4}^{\,j}\otimes
              (\mathcal O/\mathcal I_{C_4}^{T_m-2j}),
              \label{eq:quartic-nilradical-powers}\\
 \mathcal N^j/\mathcal N^{j+1}
 &\simeq\mathcal I_{S_4}^{\,j}\otimes
       \mathcal O/(\mathcal I_{S_4}+\mathcal I_{C_4}^{T_m-2j}).
              \label{eq:quartic-nilradical-layers}
\end{align}
For $2j\ge T_m$ the power is zero. If
$Q\subset\Sym^2\mathcal E$ is the image of the initial secant
equation, then the normal cone of $C_4$ in $D_m$ has algebra
\begin{equation}\label{eq:quartic-normal-cone}
 \operatorname{gr}_{\mathcal I_{C_4}}\mathcal O_{D_m}
  \simeq \Sym\mathcal E\big/
       \big(Q\cdot\Sym^{\ge T_m-2}\mathcal E\big).
\end{equation}
Its transverse equation and its complete local thickening are
constant along $C_4$ in the sense of \eqref{eq:quartic-complete-germ}.
\end{corollary}
\begin{proof}
The secant cubic is Cartier, so
$\mathcal N\simeq\mathcal I_{S_4}\otimes
\mathcal O/\mathcal I_{C_4}^{T_m-2}$. The standard associated graded
algebra of the regular embedding $C_4\subset\Pj^4$ gives
\eqref{eq:quartic-conormal-layers}. Locally,
\[
 (F^j)\cap(F\mathfrak p^{T_m-2})
      =F^j\mathfrak p^{\max\{T_m-2j,0\}}.
\]
This proves the power formula; quotienting by the next power proves
\eqref{eq:quartic-nilradical-layers}. Finally the exact homogeneous
local ideal $F\mathfrak p^{T_m-2}$ has initial ideal equal to
itself. These local descriptions glue through the intrinsic
conormal bundle and the initial secant equation, proving the
normal-cone identity.
\end{proof}
